% Blueprint content: one node per definition or numbered result of the book.
%
% Conventions used throughout:
%   \lean{...}   the Lean declaration(s) realising this node
%   \leanok      the Lean statement exists AND is proved with no `sorry`
%   \notready    stated in Lean but assumed (`sorry`), or not statable at all
% A node with no \lean{} is one today's Mathlib cannot even express; the text
% says what is missing.  Statements are original paraphrases, never the book's
% own wording.

\chapter{Foundations}

\begin{definition}[Second differential]
  \label{def:sndFDeriv}
  \lean{MorseFloer.sndFDeriv}
  \leanok
  For $f$ on a normed space, $(d^2f)_x$ is the derivative of $x \mapsto (df)_x$,
  read as a continuous bilinear map.
\end{definition}

\begin{theorem}[Transformation rule for second differentials]
  \label{thm:sndFDeriv_comp}
  \lean{MorseFloer.sndFDeriv_comp}
  \leanok
  \uses{def:sndFDeriv}
  Under a change of coordinates $\varphi$,
  \[ d^2(f\circ\varphi)_x(h,k) = d^2f_{\varphi(x)}(d\varphi\, h, d\varphi\, k)
     + df_{\varphi(x)}\bigl(d^2\varphi_x(h,k)\bigr). \]
  The correction term lies in the image of $df$, which is exactly why the second
  differential is not chart-independent in general.
\end{theorem}

\begin{corollary}[Tensoriality at a critical point]
  \label{cor:sndFDeriv_critical}
  \lean{MorseFloer.sndFDeriv_comp_of_critical}
  \leanok
  \uses{thm:sndFDeriv_comp}
  Where $df$ vanishes the correction term disappears, so $d^2f$ transforms as a
  bilinear form.  This is what makes the Hessian of a function on a manifold well
  defined at its critical points, and nowhere else.
\end{corollary}

\begin{definition}[Hessian, critical point, Morse function]
  \label{def:hessian}
  \lean{MorseFloer.hessian, MorseFloer.IsCriticalPoint, MorseFloer.IsNondegenerate,
        MorseFloer.IsMorseFunction}
  \leanok
  \uses{cor:sndFDeriv_critical}
  On a manifold, $x$ is critical for $f$ when $d f_x = 0$; the Hessian is the
  second differential of $f$ read in the preferred extended chart at $x$; the
  critical point is nondegenerate when that bilinear form is; and $f$ is a Morse
  function when all its critical points are nondegenerate.
\end{definition}

\begin{definition}[Index of a critical point]
  \label{def:index}
  \lean{MorseFloer.index, MorseFloer.morseIndex}
  \leanok
  \uses{def:hessian}
  The index is the largest dimension of a subspace on which the Hessian quadratic
  form is negative definite, i.e.\ \texttt{sigNeg} of that form.
\end{definition}

\begin{proposition}[The index is well defined]
  \label{prop:index_invariant}
  \lean{MorseFloer.index_eq_of_equivalent}
  \leanok
  \uses{def:index}
  Equivalent quadratic forms have the same index.  This is Sylvester's law of
  inertia, already available in Mathlib, and it is what makes the index
  independent of the coordinates used to compute it.
\end{proposition}

\chapter{Morse functions (Chapter 1)}

\section{\S1.1 Definition of Morse functions}

\begin{proposition}[Symmetry of the Hessian]
  \label{prop:hessian_symm}
  \lean{MorseFloer.Chapter1.hessian_symm}
  \leanok
  \uses{def:hessian}
  The Hessian is a symmetric bilinear form.  The book derives this from
  $X\cdot(Y\cdot f) - Y\cdot(X\cdot f) = [X,Y]\cdot f = df([X,Y]) = 0$; in a chart
  it is the symmetry of the second derivative of a $C^2$ function.
\end{proposition}

\begin{proposition}[Chart-independence at a critical point]
  \label{prop:hessian_chart}
  \lean{MorseFloer.Chapter1.hessian_chart_independent}
  \leanok
  \uses{cor:sndFDeriv_critical}
  Remark 1.1.2 and Exercise 1: the Hessian at a critical point does not depend on
  the chart used to compute it.
\end{proposition}

\section{\S1.2 Existence and multitude of Morse functions}

\begin{definition}[Squared distance to a point]
  \label{def:distSq}
  \lean{MorseFloer.Chapter1.distSq}
  \leanok
  The function $f_p(x) = \lVert x - p\rVert^2$ of \S1.2.a.
\end{definition}

\begin{proposition}[Differential of the squared distance]
  \label{prop:fderiv_distSq}
  \lean{MorseFloer.Chapter1.hasFDerivAt_distSq, MorseFloer.Chapter1.fderiv_distSq,
        MorseFloer.Chapter1.isCriticalPt_distSq_iff}
  \leanok
  \uses{def:distSq}
  $(df_p)_x(\xi) = 2\langle x-p, \xi\rangle$.  On a submanifold $V$ this says $x$
  is critical for $f_p$ exactly when $x-p \perp T_xV$, the characterisation the
  book feeds to Sard's theorem through the normal bundle.
\end{proposition}

\begin{proposition}[Proposition 1.2.1]
  \label{prop:1.2.1}
  \lean{MorseFloer.Chapter1.ae_isMorseFunction_distSq}
  \notready
  \uses{prop:fderiv_distSq}
  For almost every $p \in \mathbb{R}^n$, the function $f_p$ restricted to a
  submanifold $V \subseteq \mathbb{R}^n$ is a Morse function.
  \emph{Assumed:} the proof applies Sard's theorem to the endpoint map of the
  normal bundle, and Sard's theorem for manifolds is not in Mathlib.
\end{proposition}

\begin{lemma}[Lemma 1.2.2]
  \label{lem:1.2.2}
  \notready
  The normal bundle $N = \{(x,v) \in V\times\mathbb{R}^n : v \perp T_xV\}$ is a
  submanifold of $V\times\mathbb{R}^n$, and $p = x+v$ is a critical value of the
  endpoint map $E(x,v) = x+v$ exactly when the second-derivative matrix of $f_p$
  fails to be invertible.
  \emph{Not statable:} Mathlib has no normal bundle of a submanifold and no
  general submanifold type to carry it.
\end{lemma}

\begin{proposition}[Proposition 1.2.4]
  \label{prop:1.2.4}
  \notready
  Any $C^\infty$ function on a manifold embeddable in $\mathbb{R}^n$ can be
  approximated, uniformly on compact sets together with its derivatives of order
  $\le k$, by Morse functions.
  \emph{Not statable:} needs the $C^k$ topology on $C^\infty(V;\mathbb{R})$,
  which Mathlib does not define.
\end{proposition}

\begin{theorem}[Theorem 1.2.5]
  \label{thm:1.2.5}
  \notready
  \uses{prop:1.2.4}
  On a compact manifold the Morse functions form an open dense subset of
  $C^\infty(V;\mathbb{R})$.
  \emph{Not statable:} same missing topology as Proposition 1.2.4.
\end{theorem}

\section{\S1.3 The Morse lemma, index of a critical point}

\begin{theorem}[Theorem 1.3.1, the Morse lemma]
  \label{thm:1.3.1}
  \lean{MorseFloer.Chapter1.morse_lemma}
  \notready
  \uses{def:hessian}
  Near a nondegenerate critical point $c$ there are coordinates in which $f$
  equals its quadratic model: a chart $\varphi$ centred at $c$ with
  $f(y) = f(c) + \tfrac12 d^2f_c(\varphi y, \varphi y)$.  Diagonalising by
  Sylvester's law gives the book's normal form
  $f(c) - \sum_{j\le i} x_j^2 + \sum_{j>i} x_j^2$.
  \emph{Assumed:} the book's proof is an induction on dimension via the implicit
  function theorem; not yet carried out in Lean.
\end{theorem}

\begin{corollary}[Corollary 1.3.2]
  \label{cor:1.3.2}
  \lean{MorseFloer.Chapter1.isolated_of_nondegenerate}
  \leanok
  \uses{def:hessian}
  Nondegenerate critical points are isolated.  Proved as Exercise 2 asks,
  \emph{without} the Morse lemma: $df$ has an invertible strict derivative at
  $c$, so the inverse function theorem makes it a local homeomorphism and $c$ is
  its only nearby zero.
\end{corollary}

\begin{proposition}[Remark 1.3.3]
  \label{prop:1.3.3}
  \lean{MorseFloer.Chapter1.index_neg_add_index}
  \leanok
  \uses{def:index}
  A critical point of index $i$ for $f$ has index $n-i$ for $-f$.
\end{proposition}

\section{\S1.4 Examples of Morse functions}

\begin{proposition}[The torus function, \S1.4.c]
  \label{prop:1.4.c}
  \lean{MorseFloer.Chapter1.cosFactor,
        MorseFloer.Chapter1.deriv_cosFactor,
        MorseFloer.Chapter1.deriv_deriv_cosFactor,
        MorseFloer.Chapter1.deriv_cosFactor_eq_zero_iff,
        MorseFloer.Chapter1.deriv_deriv_cosFactor_ne_zero,
        MorseFloer.Chapter1.deriv_deriv_cosFactor_int,
        MorseFloer.Chapter1.deriv_deriv_cosFactor_half_odd}
  \leanok
  \uses{def:index}
  The one-variable factor $t \mapsto \cos 2\pi t$ has critical points exactly the
  half-integers, all nondegenerate, with second derivative $-4\pi^2$ at the
  integers and $+4\pi^2$ at the half-odd-integers.  Summing two copies gives the
  four critical points of $\cos 2\pi x + \cos 2\pi y$ on the torus, of indices
  $0, 1, 1, 2$.
\end{proposition}

\chapter{A little differential geometry (Chapter 14)}

This appendix is where the project's Mathlib coverage is decided: the gaps
recorded here are what force the assumed results elsewhere.

\begin{definition}[Critical point of a map, \S14.2.a]
  \label{def:14.2.a}
  \lean{MorseFloer.Chapter14.HasMaximalRank, MorseFloer.Chapter14.IsCriticalPointOfMap}
  \leanok
  A point is critical for $f\colon M\to N$ when $T_xf$ fails to have maximal
  rank.  Maximal rank is recorded as ``injective or surjective'', which is
  equivalent to the book's ``corank zero'' in either dimension regime and avoids
  a case split.
\end{definition}

\begin{proposition}[The fundamental example, \S14.2.a]
  \label{prop:14.2.a}
  \lean{MorseFloer.Chapter14.isCriticalPointOfMap_iff_eq_zero}
  \leanok
  \uses{def:14.2.a}
  For a real-valued function the general definition collapses to $T_xf \equiv 0$.
  This is what reconciles the appendix's definition with the one used throughout
  Part I.
\end{proposition}

\begin{theorem}[Theorem 14.2.1, Sard]
  \label{thm:14.2.1}
  \lean{MorseFloer.Chapter14.criticalSet, MorseFloer.Chapter14.criticalValues,
        MorseFloer.Chapter14.sard}
  \notready
  \uses{def:14.2.a}
  The critical values of a smooth map form a set of measure zero.
  \emph{Assumed}, and the most consequential gap in the project: Proposition
  1.2.1 and every transversality genericity statement in the book rest on it.
\end{theorem}

\begin{definition}[Transverse subspaces, \S14.3.a]
  \label{def:14.3}
  \lean{MorseFloer.Chapter14.IsTransverse}
  \leanok
  Two subspaces are transverse when together they span: the tangent-space
  condition $T_uM + T_uN = T_uP$ of Theorem 14.3.1.
\end{definition}

\begin{theorem}[Theorem 14.3.1, linear core]
  \label{thm:14.3.1}
  \lean{MorseFloer.Chapter14.finrank_add_finrank_of_isTransverse,
        MorseFloer.Chapter14.codim_inf_of_isTransverse,
        MorseFloer.Chapter14.not_isTransverse_of_finrank_add_lt}
  \leanok
  \uses{def:14.3}
  A transverse intersection has the expected dimension, so codimensions add; and
  subspaces whose dimensions sum below the ambient dimension cannot be
  transverse (Remark 14.3.4).  The manifold statement --- that $M\cap N$ is then
  a submanifold with $T_u(M\cap N) = T_uM\cap T_uN$ --- needs a submanifold type
  Mathlib does not have.
\end{theorem}

\begin{theorem}[Theorem 14.1.1, characterisation of submanifolds]
  \label{thm:14.1.1}
  \notready
  For $V \subseteq \mathbb{R}^n$, being a $d$-dimensional submanifold is
  equivalent to being locally the zero set of a submersion, and to being locally
  the homeomorphic image of an immersion.
  \emph{Not statable:} Mathlib has no submanifold type to carry the equivalence.
\end{theorem}

\begin{proposition}[Proposition 14.3.5 and Theorem 14.3.7]
  \label{prop:14.3.5}
  \notready
  \uses{thm:14.2.1}
  If $f \pitchfork P$ then $f^{-1}(P)$ is a submanifold; and if
  $F\colon M\times S\to N$ is transverse to $P$ then $F_s$ is transverse to $P$
  for almost every parameter $s$ --- the parametric transversality argument that
  Chapters 8 and 11 use constantly.
  \emph{Not statable:} needs submanifolds, and its proof needs Sard.
\end{proposition}

\begin{theorem}[Theorems 14.3.10--14.3.13, genericity of transversality]
  \label{thm:14.3.11}
  \notready
  \uses{prop:14.3.5}
  Maps transverse to a submanifold along a compact set form a dense open subset
  of $C^\infty(M,N)$.
  \emph{Not statable:} needs the $C^k$ topology on $C^\infty(M,N)$.
\end{theorem}

\chapter{The complex of critical points (Chapter 3)}

Mathlib has no moduli space of trajectories, so the analytic input of \S3.2
cannot be stated at all.  Everything downstream of it --- the entire algebraic
skeleton of the chapter --- is formalized and proved.

\section{\S3.1 The complex}

\begin{definition}[The chain groups and the differential]
  \label{def:3.1}
  \lean{MorseFloer.Chapter3.CritSet, MorseFloer.Chapter3.Chains,
        MorseFloer.Chapter3.dLin}
  \leanok
  \uses{def:index}
  $C_k(f)$ is the free module on the critical points of index $k$, and
  $\partial_X a = \sum_b n_X(a,b)\, b$ where $n_X(a,b)$ counts, modulo $2$, the
  trajectories of $X$ from $a$ down to $b$.  The count enters as abstract data
  `cnt`, since the geometry that produces it is not available.
\end{definition}

\begin{definition}[The broken-trajectory hypothesis]
  \label{def:brokenPairs}
  \lean{MorseFloer.Chapter3.BrokenPairs}
  \leanok
  \uses{def:3.1}
  For $\operatorname{Ind} a = \operatorname{Ind} b + 2$, the once-broken
  trajectories from $a$ to $b$ cancel in pairs:
  $\sum_c n_X(a,c)\,n_X(c,b) = 0$.  This is the \emph{conclusion} of
  Theorem 3.2.6 together with the fact that a compact $1$-manifold has an even
  number of boundary points, taken as a hypothesis because its proof is
  geometric.
\end{definition}

\begin{theorem}[$\partial_X\circ\partial_X = 0$, \S3.1.b]
  \label{thm:dd0}
  \lean{MorseFloer.Chapter3.dLin_comp_dLin}
  \leanok
  \uses{def:brokenPairs}
  A complete proof from the broken-trajectory hypothesis: the coefficient of $b$
  in $\partial\partial a$ is exactly the sum that hypothesis kills.  Only the
  geometric input is assumed.
\end{theorem}

\begin{definition}[The Morse complex and its homology]
  \label{def:morseComplex}
  \lean{MorseFloer.Chapter3.morseComplex, MorseFloer.Chapter3.morseHomology,
        MorseFloer.Chapter3.cycles, MorseFloer.Chapter3.boundaries,
        MorseFloer.Chapter3.boundaries_le_cycles}
  \leanok
  \uses{thm:dd0}
  $(C_\star(f),\partial_X)$ as an object of Mathlib's category of chain
  complexes, so the whole homological-algebra library applies, together with the
  book's description of $H_k$ as $\ker\partial_k/\operatorname{im}\partial_{k+1}$.
\end{definition}

\section{\S3.2 The geometric input}

\begin{theorem}[Theorems 3.2.2, 3.2.6 and Corollary 3.2.4]
  \label{thm:3.2}
  \notready
  The compactified space of trajectories $\bar L(a,b)$ is compact; $L(a,b)$ is
  finite when the index drops by one (which is what makes $n_X(a,b)$ meaningful);
  and when the index drops by two it is a compact $1$-manifold with boundary the
  once-broken trajectories.
  \emph{Not statable:} there is no space of trajectories, no compactification
  with the topology of \S3.2.a, no manifolds with boundary or corners, and no
  Smale condition in Mathlib.  Their joint consequence is what
  \ref{def:brokenPairs} assumes.
\end{theorem}

\begin{lemma}[Lemma 3.2.11, the Morse model]
  \label{lem:3.2.11}
  \lean{MorseFloer.Chapter3.modelFun, MorseFloer.Chapter3.modelField,
        MorseFloer.Chapter3.modelFlow, MorseFloer.Chapter3.modelTransition,
        MorseFloer.Chapter3.modelFlow_log_eq_modelTransition,
        MorseFloer.Chapter3.modelFun_modelTransition}
  \leanok
  In a Morse chart the flow carries one end of the chart to the other by
  $(x^-,x^+) \mapsto (\lVert x^+\rVert/\lVert x^-\rVert\cdot x^-,\;
  \lVert x^-\rVert/\lVert x^+\rVert\cdot x^+)$, and the model function decreases
  strictly along the flow.  Proved in full.
\end{lemma}

\section{\S3.3 Integer coefficients}

\begin{proposition}[Reduction mod 2]
  \label{prop:3.3}
  \lean{MorseFloer.Chapter3.BrokenPairs.map, MorseFloer.Chapter3.BrokenPairs.intCast}
  \leanok
  \uses{def:brokenPairs}
  The same construction over $\mathbb{Z}$ using orientations; the mod $2$ counts
  are the reductions of the integral ones, so the broken-trajectory hypothesis
  descends.
\end{proposition}

\section{\S3.4 Independence of the data}

\begin{theorem}[Theorem 3.4.2, algebraic steps]
  \label{thm:3.4.2}
  \lean{MorseFloer.Chapter3.coneDiff, MorseFloer.Chapter3.chainMap_of_coneDiff_comp,
        MorseFloer.Chapter3.sub_mem_range_of_chainHomotopy}
  \leanok
  \uses{def:morseComplex}
  Two steps of the invariance proof are pure algebra and are proved: a
  differential triangular in a decomposition $C_k(f_0)\oplus C_{k+1}(f_1)$
  squares to zero exactly when its off-diagonal part is a chain map, and two
  maps differing by a chain homotopy agree on homology classes.
  The theorem itself --- that the homology does not depend on $(f,X)$ --- links
  those steps through the geometry and is not statable for abstract count data.
\end{theorem}
